\chapter{Tổng quan đề tài}
Băng tải là một hệ thống cơ khí thiết yếu, được ứng dụng rộng rãi trong các quy trình sản xuất và hoạt động xếp dỡ hàng hóa nhằm mục đích vận chuyển vật liệu \cite{ufot2017fuzzy}. Trong các hệ thống phân loại, việc kiểm soát tốc độ băng tải là yếu tố then chốt. Tốc độ ổn định là điều kiện tiên quyết để các cảm biến có thể thu thập và xử lý dữ liệu chính xác, giảm thiểu nhiễu động, đồng thời đảm bảo sự ổn định chung của dây chuyền sản xuất, cho phép công nhân kịp thời khắc phục sự cố.\\

Các hệ thống này thường được dẫn động bằng động cơ một chiều (DC). Để điều chỉnh tốc độ, bộ phận điều khiển sẽ liên tục tính toán độ lệch giữa tốc độ mong muốn và tốc độ thực tế. Tín hiệu điều khiển sau đó được tạo ra thông qua nhiều phương pháp khác nhau, bao gồm các thuật toán thông minh, Bộ điều khiển Logic Khả trình (PLC - Programmable Logic Controller) – đóng vai trò "bộ não" điều khiển các quá trình cơ khí, Bộ điều khiển PID, hoặc các thuật toán kết hợp.\\

Đã có nhiều nghiên cứu áp dụng các giải pháp này, chẳng hạn như Umoren và cộng sự \cite{umorn2016design} đã ứng dụng bộ điều khiển PID cho băng tải công nghiệp nhằm đồng bộ hóa tốc độ giữa các công đoạn trong nhà máy bia. Tương tự, Ufot và cộng sự \cite{ufot2017fuzzy} đã sử dụng Bộ điều khiển Logic Mờ (FLC) để điều khiển đồng bộ tốc độ dây chuyền. Tuy nhiên, các hệ thống hiện tại vẫn cần được cải thiện hiệu năng và giảm sai số truyền, cho phép vật liệu di chuyển nhanh hơn và đồng bộ hơn với tốc độ vận hành của toàn bộ dây chuyền.\\

Do đó, nhóm nghiên cứu tập trung vào việc thiết kế bộ điều khiển PID cho động cơ DC và triển khai trên hệ thống phân loại màu sắc. Mục tiêu là tối ưu hóa hiệu quả điều khiển, đảm bảo hệ thống đạt độ chính xác cao nhất trong quá trình phân loại.

